\documentclass[11pt]{article}

\usepackage[margin=1in]{geometry}
\usepackage{amsmath,amssymb,amsthm}
\usepackage{booktabs}
\usepackage{enumitem}
\usepackage{hyperref}

\newcommand{\ket}[1]{\lvert #1 \rangle}
\newcommand{\Z}{\mathbb{Z}}
\newcommand{\R}{\mathbb{R}}
\newcommand{\Exp}{\mathbb{E}}

\theoremstyle{plain}
\newtheorem{fact}{Key fact}
\theoremstyle{definition}
\newtheorem*{defn}{Definition}

\title{\textbf{Breaking the Silent Group}\\ \large quantum preliminaries, the task, and how the solver works}
\author{Raiden Shor-gun --- QFF Hackathon, Track 2}
\date{}

\begin{document}
\maketitle

\section{The problem in one paragraph}

A sealed quantum device holds one element of a finite abelian group $G$
($|G| \le 2^{60}$, structure unknown) in a register we can never look at.
The only operation is a controlled multiplication: hand the device a
control qubit and an instruction ``multiply by $h^k$'' and, conditioned on
the control being $\ket{1}$, it applies $U_{h^k}$ to its register.
Three parties run a Diffie--Hellman-style exchange over this group:
public broadcasts $g, A=g^a, B=g^b, C=g^c$, shared key $K=g^{abc}$.
Given $(g,A,B,C)$ we must report
\[
  \mathrm{order} = \operatorname{ord}(g), \qquad
  \mathrm{secret} = abc \bmod \operatorname{ord}(g),
\]
or \texttt{"none"} if any broadcast is \emph{forged} --- an element that is
not a power of $g$ at all. Every oracle call is billed
$\max(1,\operatorname{bitlength}(k))$ \emph{queries} from an 80{,}000-query
budget, and every call adds Gaussian phase noise of standard deviation
$\sigma\sqrt{\mathrm{cost}}$ turns, $\sigma = 0.002$.

Classically this is order finding plus three discrete logs --- the exact
pair of problems Shor's algorithm solves. The twist is that the group is
a black box: the only thing the hardware ever reveals is \emph{phase}.

\section{Quantum preliminaries (the only part we need)}

\subsection{Qubits, phases, and interference}

A qubit is a unit vector $\alpha\ket{0}+\beta\ket{1}$; measuring returns a
bit with probabilities $|\alpha|^2, |\beta|^2$. The two gates we use are
\[
  H = \tfrac{1}{\sqrt2}\begin{pmatrix}1&1\\1&-1\end{pmatrix},
  \qquad
  P(\theta) = \begin{pmatrix}1&0\\0&e^{i\theta}\end{pmatrix}.
\]
$H$ creates the uniform superposition $\ket{+}=(\ket0+\ket1)/\sqrt2$.
A phase on the $\ket1$ branch is invisible by itself, but a second $H$
lets it interfere:
\[
  H\,P(2\pi\theta)\,H\ket0
  = \tfrac{1+e^{2\pi i\theta}}{2}\ket0 + \tfrac{1-e^{2\pi i\theta}}{2}\ket1,
  \qquad \Pr[1] = \sin^2(\pi\theta).
\]
So $H$--phase--$H$--measure is a \emph{phase detector}: it reads out
$\theta$ with bias $\sin^2(\pi\theta)$, deterministically when
$\theta\in\{0,\tfrac12\}$.

\subsection{Phase kickback}

If $U\ket{u}=e^{i\theta}\ket{u}$ (i.e.\ $\ket{u}$ is an eigenvector), then
applying controlled-$U$ to $\ket{+}\ket{u}$ gives
\[
  \frac{\ket0 + e^{i\theta}\ket1}{\sqrt2}\otimes\ket{u}.
\]
The target register is \emph{unchanged}; the eigenphase moves onto the
control's $\ket1$ branch. This is phase kickback, and it is the entire
mechanism by which we read the device: our one qubit acts as the control,
the group register is the eigenvector, and $\theta$ carries information
about $h$.

\subsection{Characters of a finite abelian group}

A \emph{character} is a homomorphism $\chi : G \to \R/\Z$, i.e.\
$\chi(xy)=\chi(x)+\chi(y) \pmod 1$. For $G=\Z_{d_1}\times\cdots\times\Z_{d_m}$
they are exactly
\[
  \chi_c(h) = \sum_i \frac{c_i h_i}{d_i} \pmod 1, \qquad c\in G,
\]
so $G$ has $|G|$ characters, indexed by $G$ itself. Two consequences run
the whole algorithm:

\begin{itemize}[nosep]
  \item On the cyclic subgroup $\langle g\rangle$ of order $n$,
        every character has $\chi(g)=j/n$ for some
        $j\in\{0,\dots,n-1\}$. \textbf{The denominator of $\chi(g)$ is the
        order we are asked for.}
  \item $\chi$ is a homomorphism, so $\chi(g^a)=a\,\chi(g)$. Within one
        shot the same $\chi$ applies to every element, hence
        $\chi(A)=a j/n$: \textbf{the secret exponents appear as linear
        congruences between measured phases.}
\end{itemize}

\subsection{Why the device is a ``random character generator''}

All multiplications $U_h$ commute ($G$ abelian), so they share an
eigenbasis: one eigenvector $\ket{\hat\chi}$ per character, with
$U_h\ket{\hat\chi} = e^{2\pi i\,\chi(h)}\ket{\hat\chi}$.
The register starts in the identity state, which in this basis is the
uniform superposition
\[
  \ket{e} = \frac{1}{\sqrt{|G|}}\sum_\chi \ket{\hat\chi}.
\]

\begin{fact}
The register is never measured, so the $\chi$-branches never interfere.
Each shot of a circuit therefore behaves \emph{exactly} as if
(i)~a character $\chi$ is drawn uniformly at random, fixed for the whole
shot, and (ii)~each \texttt{oracle\_gate(h,k)} is a phase gate
$P\!\big(2\pi k\,\chi(h)\big)$ on our qubit.
\end{fact}

This is not an approximation --- the simulator literally does this
(sampling $\chi$ per shot and multiplying the $\ket1$ amplitude by the
phase), and it is exact at any group size. The ``quantum computer'' in
this task is a device that turns black-box group multiplication into
random, mutually consistent phase samples. Everything after this point is
signal processing and number theory.

\section{Reading a phase with one qubit}

\subsection{Iterative phase estimation = semiclassical QFT}

An oracle call \texttt{oracle\_gate(h,}$2^k$\texttt{)} implements
$P(2\pi\cdot 2^k\chi(h))$. Write the unknown phase in binary,
$\chi(h)=0.b_1b_2b_3\dots$; then
\[
  2^k\chi(h) \bmod 1 = 0.b_{k+1}b_{k+2}\dots
\]
--- each power left-shifts the binary point past the bits already read.
We therefore read bits from the bottom up: power $2^{m-1}$ isolates
$0.b_m$ (clean measurement), and before reading bit $b_k$ we \emph{undo}
the contribution of the lower bits we already measured,
\[
  \text{if bit } b_{k+l}=1:\quad P\!\big(-2\pi/2^{l+1}\big),
\]
implemented with mid-circuit measurement, \texttt{reset}, and
\texttt{if\_test}. This is the inverse quantum Fourier transform with its
controlled rotations replaced by classical feed-forward
(Griffiths--Niu semiclassical QFT). The register is never touched.

Two consequences:

\begin{itemize}[nosep]
  \item \textbf{Every circuit uses 1 qubit --- the minimum possible.}
        A control channel is unavoidable, so 0 is impossible and
        1 suffices.
  \item \textbf{Cost.} Step $k$ uses power $2^k$, billed $k+1$ queries,
        so an $m$-bit estimate costs
        $C(m)=\sum_{k=0}^{m-1}(k+1)=m(m+1)/2$ queries --- quadratic in
        precision, and optimal: one measurement yields at most one bit,
        and resolving scale $2^{-k}$ requires a power $\ge 2^{k-1}$.
\end{itemize}

\subsection{Noise}

Each oracle call adds an independent Gaussian error
$\varepsilon\sim\mathcal N(0,\sigma^2\cdot\mathrm{cost})$ \emph{turns} to
the kicked phase. At the step using power $2^k$ the phase detector is
binary, so the bit flips with probability
\[
  q_k = \Exp\!\left[\sin^2(\pi\varepsilon)\right]
      = \frac{1-e^{-2\pi^2\sigma^2(k+1)}}{2}
      \approx \pi^2\sigma^2(k+1).
\]
The expensive steps are the noisy steps. A flipped bit at weight $2^p$
shifts the estimate by exactly $\pm 2^p$ LSBs, so \textbf{list decoding}
--- try undoing each single $\pm 2^p$ shift and re-fit --- recovers most
corrupted shots. At competition noise a 132-bit shot has
$\lambda=\sum q_k\approx 0.32$: 72\% arrive clean, $\sim$96\% are usable
after decoding one flip.

\section{The algorithms}

\subsection{Identity testing: the workhorse}

To test whether $t=\prod h_i^{k_i}$ is the identity, run
$H$, $\prod$\texttt{oracle\_gate}$(h_i,k_i)$, $H$, measure. Then
$\Pr[1]=\sin^2(\pi\chi(t))$, and:

\begin{itemize}[nosep]
  \item $t=e$: $\chi(t)=0$ for every $\chi$ $\Rightarrow$ \emph{always} 0.
  \item $t\ne e$ of order $s$: $\chi(t)$ is uniform over
        $\{0,1/s,\dots,(s-1)/s\}$ $\Rightarrow$ $\Pr[1]=\tfrac12$
        \emph{exactly}, independently of everything else.
\end{itemize}

So the test is a perfect unbiased coin \emph{unless} $t=e$. Sixteen
consecutive zeros accept with false-accept probability $\approx2^{-16}$;
three ones reject. Every verification in the solver --- order
confirmation, log verification, forgery proofs, relation certification
--- reduces to this $\chi$-independent test.

\subsection{Stage 1: the order of $g$}

Estimate $\chi(g)\approx M/2^m$ to $m=2\cdot\mathrm{bound}+12$ bits. The
true value is $j/n$ with $n\le 2^{60}$, so we take the \emph{simplest
fraction in the $\pm8$-LSB window} (Stern--Brocot descent): the
lowest-denominator fraction near the estimate, whose denominator is the
candidate order. The window is so narrow that valid fractions cover only
$\sim$0.1\% of $[0,1)$ --- a corrupted shot almost never fakes a fit, and
a real one always has $j/n$ waiting inside it. Every candidate is then
confirmed by testing $g^c\stackrel?=e$.

Two refinements matter for cost:

\begin{itemize}[nosep]
  \item \textbf{Probe shot.} Sizing every shot for the worst case
        ($n\approx2^{60}$, 132 bits, 8.8k queries) wastes 93\% of the
        budget when $n$ is small. We first shoot 84 bits, which covers
        $n\le2^{36}$ for 3.6k queries.
  \item \textbf{Lattice reuse.} If the probe fails we shoot 116 bits.
        Neither shot alone covers $n\le2^{60}$, but the two estimates
        $e_1/2^{m_1}, e_2/2^{m_2}$ share denominator $n$ --- recovering it
        is simultaneous Diophantine approximation: find small $q$ with
        both $q e_i/2^{m_i}$ within $8q$ of integers, i.e.\ a short
        vector of the lattice spanned by
        $(8,e_1,e_2)$, $(0,2^{m_1},0)$, $(0,0,2^{m_2})$.
        A 3-dimensional LLL reduction finds it. The failed probe is not
        wasted, and the answer is still confirmed by $g^n=e$, so the
        lattice makes us cheaper, never riskier.
\end{itemize}

If a confirmed $c$ is a multiple of $\operatorname{ord}(g)$, the excess
factors divide $\gcd(j,c)$ and are stripped later by identity tests on
$g^{c/p^t}$.

\subsection{Stage 2: discrete logs}

One joint shot now estimates $\chi(g),\chi(A),\chi(B),\chi(C)$
\emph{with the same $\chi$}. Rounding each estimate to the $(1/n)\Z$
grid gives $j/n$ and $r_X/n$ with
\[
  a\,j \equiv r_A \pmod n
\]
(and similarly for $b,c$) --- a linear congruence whose solutions are a
class $x_0 + tM$, $t<n/M$. Each candidate is verified by the identity
test $X\cdot g^{-x}\stackrel?=e$; consistent rounds are merged across
shots by CRT, and large classes are tested wholesale via
$(Xg^{-x_0})^{n/M}\stackrel?=e$.

\subsection{Stage 3: forgery --- proving a negative}

A forged $X\notin\langle g\rangle$ comes in two grades:

\begin{itemize}[nosep]
  \item \textbf{Easy:} $X^n\ne e$. One identity test proves it.
  \item \textbf{Hard:} $X\in G[n]$ (the $n$-torsion) but
        $X\notin\langle g\rangle$ --- e.g.\ \texttt{public\_2}'s $B$.
        Every round looks clean, yet every candidate fails verification
        or the congruence is inconsistent ($r\notin j\Z_n$).
\end{itemize}

We count \emph{clean rounds that contradict $\langle g\rangle$} and call a
forgery past a threshold. The subtlety is that \emph{honest} elements
occasionally produce clean-but-wrong rounds: a noise flip at the step
with power $2^k$ shifts $\chi\cdot n$ by $n/2^{k+1}$, which goes
unnoticed precisely when that is near an integer --- worst case
$n\approx 2^k$. The honest-bad-round rate $p(n)$ is therefore computed
per-instance from the flip model (matching simulation within $\sim$20\%),
and the threshold is the smallest $e\ge3$ with $(3p)^e\le10^{-5}$.
A forgery in a tiny factor like $\Z_2$ stays invisible until
$\gcd(j,n)$ is coprime to it --- about half the rounds --- which is the
main residual cost.

\subsection{Graceful degradation}

When the budget or the noise wins, the solver still returns the best
order found with \texttt{"none"}: the order alone is worth 0.3 points,
and a forgery is the usual reason the logs cannot finish. If the logs
stage ever disproves the order ($g^n\ne e$ caught by a strict backstop
test), the order stage restarts rather than reporting a divisor ---
a wrong order scores 0.

\section{Open track: the full structure of $\langle g,A,B,C\rangle$}

Each shot's $\chi$ restricted to $H=\langle g,A,B,C\rangle$ is a
\emph{uniform character of $H$}. One joint phase-estimation shot on the
four generators reads $\chi(\mathrm{gens})=r/E$ with $r\in\Z_E^4$
($E$ the exponent of $H$). The vectors $r$ span a subgroup
$R\le\Z_E^4$ with $R\cong\widehat H\cong H$. Then:

\begin{itemize}[nosep]
  \item the \emph{relation lattice}
        $\Lambda=\{x\in\Z^4 : g^{x_0}A^{x_1}B^{x_2}C^{x_3}=e\}$
        is the scaled dual $E\,(R+E\Z^4)^{-1}$;
  \item $H\cong\Z^4/\Lambda$, and its Smith normal form gives the
        invariant factors $H\cong\Z_{n_1}\times\Z_{n_2}\times\cdots$;
  \item $H/\langle g\rangle\cong\Z^3/\mathrm{proj}(\Lambda)$ is trivial
        iff nothing is forged, and its size measures how far forgeries
        sit outside $\langle g\rangle$.
\end{itemize}

Nothing is taken on trust: every basis relation of $\Lambda$ is
re-verified on the device with the identity test, and a corrupted shot
that inflates $R$ is caught by leave-one-out (dropping it shrinks $R$
while every recovered relation verifies). On \texttt{public\_2} this
recovers $H\cong\Z_{585585}\times\Z_{37199287125}$ with
$H/\langle g\rangle\cong\Z_{585585}$ --- the forged $B$ lives one factor
outside $\langle g\rangle$.

\section{What failed on the way}

\begin{itemize}[nosep]
  \item \textbf{Worst-case order shot.} Always paying for $n\le2^{60}$
        made the order shot 93\% of cost on small groups. Found by the
        cost model; fixed by the probe + lattice.
  \item \textbf{Smooth $n$ + forgery.} A $\Z_5$ factor forgery with a
        smooth $n$ cost 60--80k queries every run: $\gcd(j,n)$ was large,
        so no round ever pinned a candidate class. Found by randomized
        stress; fixed by testing the whole class at once.
  \item \textbf{Rare false accepts} ($\sim$1/3000 runs): a wrong log, or
        an order silently shrunk to a divisor (0 points). Found only by
        10k-run stress sweeps; fixed with strict tests (18 zeros) where
        candidates are likely wrong, plus the $g^n\stackrel?=e$
        backstop that restarts the order stage.
  \item \textbf{Inherent remaining cost:} very smooth $n$ (median 33k)
        and forgeries hidden in tiny torsion factors (median 47k) are
        properties of per-round phase estimation, not fixable bugs.
\end{itemize}

\section{Results}

\begin{center}
\begin{tabular}{lcc}
\toprule
 & v1 & final \\
\midrule
official check & 20/20 correct & 20/20 correct \\
\texttt{public\_1} (honest, $n\sim2^{56}$) & 1.33 & 1.33 \\
\texttt{public\_2} (forged, $n\sim2^{35}$) & 1.31 & 1.38 \\
random stress (4{,}500 runs) & 1.39 & 1.41, 0 wrong \\
qubits per circuit & 1 & 1 \\
\bottomrule
\end{tabular}
\end{center}

Solver degrades gracefully: correct up to $\sim$1.5$\times$ competition
noise, and when it fails it usually still returns the right order.

\section*{The one-sentence version}

\emph{Each shot draws a random character $\chi$; kickback turns group
multiplication into phase gates $k\cdot\chi(h)$; a one-qubit semiclassical
QFT reads those phases bit by bit; and denominators, congruences and
identity tests turn phases into the order, the logs, and a certified
verdict on forgery.}

\end{document}
